\RequirePackage{mymoniker}    % Replace CarrollCD in mymoniker.sty with your LastnameFirstNameMiddleName
\RequirePackage{dissertation} % required extra stuff
\documentclass{scrartcl}

\usepackage{color,hyperref}
% Uses hyperref to link DOI
\newcommand\doilink[1]{\href{http://dx.doi.org/#1}{#1}}
\newcommand\doi[1]{doi:\doilink{#1}}

\usepackage{dirtytalk}
% Packages
\usepackage[utf8]{inputenc}
\usepackage[backend=bibtex, style=authoryear]{biblatex} % You can change the style if you prefer a different citation style
\addbibresource{References2.bib}


\begin{document}

%\begin{abstract}
%  \input{\abstractpath/Abstract\mymoniker}
%\end{abstract}

\section{Teaching Philosophy}

\par I believe teaching is an individualized journey for each student. It is my belief that my job is centered around aiding a student convince themselves that they have transitioned from a state of confusion to understanding. Each person has different innate abilities to convince themselves that they've understood what they didn't before, as well as different sensitivities to being in a state of confusion. As a teacher, I understand that the most helpful that I can be to a learner is supporting them, encouraging them, and making sure that I've convinced myself of the same material they are trying to learn -- these are all necessary to establish the proper trust between teacher and student.

\par My view on teaching was born mainly out of my experiences as a student. I've been a student in academic environments for nearly my entire life. As I've progressed through primary and secondary education, I viewed learning as fun --my most fond memories of those times were of my mathematics teachers in high school and my economics teacher in my final year.

\par Unsurprisingly, I went on to study economics and mathematics as an undergraduate. However, my relationship with learning changed from a fun experience to a much more competitive one. It was during these times I began to realize that, with the right amount of effort and repetition, no topic would be impossible to learn. Said differently, any topic would be impossible to learn without giving the discipline the proper amount of focus and attention it requires. I had many mentors during this time, both inside and outside of the economics department. They openeed my eyes to the possibilites of learning and producing knowledge as a profession.

\par I was sold. I wanted to earn a PhD in economics, and so I pursued it. After many research assistantships, both before and during my time in graduate school, my relationship with learning changed again. But this time, I believe, it became a more mature perspective. I went from comparing myself to my peers, through grades and research ideas, to assesing my understanding of conccepts as it progressed over time. One of the pivotal realizations I had was that some material would be more fresh on my mind than others. This taught me that preparation is key and that as long as I have a plan and prepare well, I am generally able to effectively communicate any of the concepts and ideas I've learned over the years. I also started to notice that it would take me less and less preparation time and that I could even teach myself related, but essentially new topics very quickly. This has been very useful in the many experiences I've had as a teaching assistant and private tutor while complete my dissertation.

\par So where am I now in my own learning journey? Although one never truly stops learning, I think I am in a unique position to not only properly convey the concepts I've come across in economics and mathematics effectively to new learners, but also to guide others in their own personal journeys. At each step of the way, I've had others supporting me in my endeavors, and I am confident that I can play a similar role for others.

\section{Teaching Experience}

\par As a graduate student, I've been a teaching assistant (TA) mainly for macroeconomics courses. These courses include introductory and intermediate level macroeconomics, as well as a seminar course in macroeconomic models and policy. 

\par I've also done a fair amount of tutoring in my academic career. I've tutored introductory and intermediate courses in mathematics and economics during my undergraduate years. During graduate school, I've also done private tutoring as an independent contractor. These included a role as a research mentor for high school students, as well as weekly tutoring for AP and college level macroeconomics, microeconomics, and calculus at various levels.

\par These experiences not only kept me up to speed in my own studies, but also exposed me to the various opportunities to teach the concepts I've learned over the years. In my view, being able to support other budding learners is my way of paying it forward since I've received similar help in my own studies.

\section{Preferred Courses}

\par To conclude, I am confident that I could teach introductory and intermediate macroeconomics and microeconomics courses at the college level. These concepts are mostly fresh on my mind, and the ones that aren't can be recovered swiftly. Although tutoring has not been my main priority during my graduate studies, it has been a nice experience that I would continue to do if I had the free time after graduation. This has led me to seriously consider pursuing opportunities to be an instructor in the future.

\par


\end{document}